\begin{tabularx}{\textwidth}{L{35mm}X X}
\toprule
\textbf{ER szerkezet} & \textbf{Relációs leképezés} & \textbf{Megjegyzés} \\
\midrule
Normál egyed & Külön reláció lesz belőle. & Az egyed tulajdonságai mezőkké válnak. \\
Kulcstulajdonság & Elsődleges kulcs mező vagy mezőcsoport. & Az egyedintegritás alapja. \\
Atomi tulajdonság & Egyszerű mező a relációban. & A mező típusa és megkötése a domainből és üzleti szabályból adódik. \\
Összetett tulajdonság & Csak az elemi résztulajdonságok kerülnek mezőként a relációba. & Például cím helyett város, utca, házszám. \\
Többértékű tulajdonság & Külön reláció, amely FK-val hivatkozik az eredeti egyedre. & Megakadályozza, hogy egy mezőben listát tároljunk. \\
1:N kapcsolat & Idegen kulcs az N oldali, gyerek relációban. & Kötelező kapcsolatnál az FK általában nem lehet NULL. \\
1:1 kapcsolat & Idegen kulcs valamelyik oldalon, gyakran az opcionális vagy függő oldalon. & Egyedi feltétel is kellhet az FK-ra. \\
N:M kapcsolat & Külön kapcsolótábla két idegen kulccsal. & A kapcsolattulajdonságok is ebbe a táblába kerülnek. \\
Gyenge egyed & Külön reláció; a tulajdonos FK-ja az összetett PK része lesz. & Az azonosítás a meghatározó kapcsolaton keresztül történik. \\
\bottomrule
\end{tabularx}
